% Gemini theme
% https://github.com/anishathalye/gemini
%
% We try to keep this Overleaf template in sync with the canonical source on
% GitHub, but it's recommended that you obtain the template directly from
% GitHub to ensure that you are using the latest version.

\documentclass[final]{beamer}

% ====================
% Packages
% ====================

\usepackage[T1]{fontenc}
\usepackage{lmodern}
\usepackage[size=custom,width=60.96,height=91.44,scale=0.95]{beamerposter}
\usetheme{gemini}
\usecolortheme{gemini}
\usepackage{graphicx}
\usepackage{booktabs}
\usepackage{tikz}
\usepackage{pgfplots}
\pgfplotsset{compat=1.14}

% ====================
% Lengths
% ====================

% If you have N columns, choose \sepwidth and \colwidth such that
% (N+1)*\sepwidth + N*\colwidth = \paperwidth
\newlength{\sepwidth}
\newlength{\colwidth}
\setlength{\sepwidth}{0.025\paperwidth}
\setlength{\colwidth}{0.3\paperwidth}

\newcommand{\separatorcolumn}{\begin{column}{\sepwidth}\end{column}}

% ====================
% Title
% ====================

\title{Linguistic MDL Explanations with LLMs \\ \large via Occam-Style Prompt Discovery}

\author{Yassine Meliani \inst{1} \and George Kontorousis \inst{1} \and Eric Elmonzino \inst{2}}

\institute[shortinst]{\inst{1} McGill University \samelineand \inst{2} MILA}

% ====================
% Footer (optional)
% ====================

% \footercontent{
%   \href{https://www.example.com}{https://www.example.com} \hfill
%   ABC Conference 2025, New York --- XYZ-1234 \hfill
%   \href{mailto:alyssa.p.hacker@example.com}{alyssa.p.hacker@example.com}}
% (can be left out to remove footer)

% ====================
% Logo (optional)
% ====================

% use this to include logos on the left and/or right side of the header:
\logoright{
\includegraphics[height=7cm]{resources/McGill_University_CoA.png}}

% ====================
% Body
% ====================

\begin{document}

\begin{frame}[t]
\begin{columns}[t]
\separatorcolumn

\begin{column}{\colwidth}

  \begin{alertblock}{Overview}
    We use a Minimum Description Length (MDL) framework, grounded in Occam's razor, 
    to discover a natural-language explanation $z$ for observed data $x$ that is both 
    concise and predictive. A frozen LLM scores candidate explanations, and this signal
    is used to train an explanation policy via reinforcement learning. 
    Our goal is an amortized system that can generate useful explanations for new datasets 
    without hand-crafted prompts.
  \end{alertblock}

  % \begin{block}{Motivation}

  %   If prompts are viewed as natural-language models, then
  %   explanation discovery should not be treated as prompt
  %   tinkering alone. We want a system that can automatically
  %   identify explanations that capture the underlying pattern
  %   in the data rather than merely restating the observations.
  
  %   % \begin{figure}
  %   %   \centering
  %   %   \begin{tikzpicture}[scale=6]
  %   %     \draw[step=0.25cm,color=gray] (-1,-1) grid (1,1);
  %   %     \draw (1,0) -- (0.2,0.2) -- (0,1) -- (-0.2,0.2) -- (-1,0)
  %   %       -- (-0.2,-0.2) -- (0,-1) -- (0.2,-0.2) -- cycle;
  %   %   \end{tikzpicture}
  %   %   \caption{A figure caption.}
  %   % \end{figure}
  
  %   This requires balancing two competing goals: explanations
  %   should be simple enough to remain interpretable, but strong
  %   enough to predict the data well. That trade-off motivates
  %   our use of the Minimum Description Length principle.
  
  % \end{block}

  \begin{block}{Motivation}

\begin{center}
\includegraphics[width=0.45\linewidth]{poster/resources/plot_good_fit.png}
\hfill
\includegraphics[width=0.45\linewidth]{poster/resources/plot_overfit.png}
\end{center}

\smallskip

\textbf{Which model is a better explanation of the data?}

\smallskip

Most people prefer the simple quadratic fit (left) over the
highly complex curve (right), even though both can match the data.

\smallskip

\textbf{Why?} We naturally apply \emph{Occam’s razor}:
prefer explanations that are both simple and predictive.

\smallskip

This intuition is captured by the Minimum Description Length (MDL)
principle:

\begin{equation*}
\mathcal{L}(z, x) = -\log p(z) - \log p(x \mid z)
\end{equation*}

\begin{itemize}
  \item $-\log p(z)$: complexity of the explanation
  \item $-\log p(x \mid z)$: how well it fits the data
\end{itemize}

\textbf{Goal:} find explanations that balance simplicity and accuracy.

\end{block}

  \begin{block}{Minimum Description Length Framework}

    We model explanation discovery as selecting a natural-language
    hypothesis $z$ that gives the shortest description of the
    observed data $x$. The total score is a two-part code:
  
    \begin{equation*}
      \mathcal{L}_{\mathrm{MDL}}(z, x)
      =
      \underbrace{-\log p^*(z)}_{\text{cost of describing } z}
      +
      \underbrace{-\log p^*(x \mid z)}_{\text{cost of describing } x \text{ given } z}
    \end{equation*}

    \begin{itemize}
      \item Short explanations get a lower prior cost.
      \item Predictive explanations get a lower conditional data cost.
      \item The best explanation is the one that minimizes the
      total description length.
    \end{itemize}
  
  \end{block}

  % \begin{block}{Method}

  %   Starting from a pretrained language model, we instantiate
  %   two copies: a frozen evaluator $p^*$ and a trainable policy
  %   $\pi_\theta(z \mid x)$ obtained by LoRA adaptation. Given
  %   observed data $x$, the policy samples a candidate explanation
  %   $z$, and the frozen evaluator scores it with an MDL-style reward
  
  %   \begin{equation*}
  %     R(z, x) = \lambda \log p^*(z) + \log p^*(x \mid z),
  %   \end{equation*}
  
  %   where the first term favors concise explanations under the
  %   model prior and the second favors explanations that make the
  %   observed data likely. In our implementation, $p^*(x \mid z)$
  %   is computed by conditioning on the candidate explanation and
  %   evaluating the likelihood of the observed target sequence.
  %   The policy is optimized with REINFORCE:
  %   \begin{equation*}
  %     \mathcal{L}_{\mathrm{policy}}
  %     = -(R-b)\log \pi_\theta(z \mid x),
  %   \end{equation*}
  %   where $b$ is a moving-average baseline used to reduce gradient
  %   variance.
  
  %   The key idea is to learn an \emph{amortized} explanation policy
  %   while keeping the scoring model fixed. This preserves a stable
  %   MDL objective, reduces reward drift, and contrasts with
  %   per-dataset hypothesis search.
  % \end{block}
\end{column}

\separatorcolumn


\begin{column}{\colwidth}

\begin{block}{Method}

  We use a pretrained LLM in two roles:
  a frozen evaluator $p^*$ and a trainable policy
  $\pi_\theta(z \mid x)$ (via LoRA).

  Given data $x$, the policy samples an explanation $z$,
  which is scored by an MDL-style reward:
  \begin{equation*}
    R(z, x) = \lambda \log p^*(z) + \log p^*(x \mid z).
  \end{equation*}

  The first term favors concise explanations,
  while the second favors explanations that predict the data.

  The policy is optimized with REINFORCE:
  \begin{equation*}
    \mathcal{L}
    = -(R - b)\log \pi_\theta(z \mid x),
  \end{equation*}
  where $b$ is a baseline.

  \smallskip

  \textbf{Key idea:}
  learn a policy that proposes explanations maximizing
  simplicity and predictive accuracy under a fixed evaluator.

\end{block}
    \begin{exampleblock}{Prompt Templates}

  \vspace{0.5em}

  \textbf{Prior term: $\boldsymbol{-\log p(z)}$}
  

  \noindent
  \texttt{System: You are a careful pattern-discovery assistant.}
  
  \noindent
  \texttt{Write a concise natural-language rule for a number sequence.}
  
  \noindent
  \texttt{Return only the rule.}
  

  \noindent
  \texttt{User: Write a plausible concise natural-language rule}
  
  \noindent
  \texttt{for a number sequence.}
  

  \noindent
  \texttt{Assistant: \{z\}}



  \textbf{Likelihood term: $\boldsymbol{-\log p(x \mid z)}$}
  
  \smallskip
  \noindent
  \texttt{System: You are a careful sequence predictor. Given a rule,}
  
  \noindent
  \texttt{generate the number sequence implied by that rule.}
  

  \noindent
  \texttt{User: Rule: \{z\}}
  

  \noindent
  \texttt{Output EXACTLY the first 10 terms as comma-separated}
  
  \noindent
  \texttt{integers, no spaces, no text.}
  

  \noindent
  \texttt{Assistant: \{x\}}


  \textbf{Conditional baseline: $\boldsymbol{-\log p(z \mid x)}$}
  

  \noindent
  \texttt{System: You are a careful pattern-discovery assistant.}
  
  \noindent
  \texttt{Given an observed number sequence, propose a concise rule}
  
  \noindent
  \texttt{that generated it.}
  
  \noindent
  \texttt{User: Observed sequence: 1, 1, 2, 5, 12, 27, 58, 121, 248, 503}
  

  \noindent
  \texttt{What is the most likely rule?}
  

  \noindent
  \texttt{Assistant: \{z\}}

\end{exampleblock}

  % \begin{alertblock}{A highlighted block}

  %   This block catches your eye, so \textbf{important stuff} should probably go
  %   here.

  %   Curabitur eu libero vehicula, cursus est fringilla, luctus est. Morbi
  %   consectetur mauris quam, at finibus elit auctor ac. Aliquam erat volutpat.
  %   Aenean at nisl ut ex ullamcorper eleifend et eu augue. Aenean quis velit
  %   tristique odio convallis ultrices a ac odio.

  %   \begin{itemize}
  %     \item \textbf{Fusce dapibus tellus} vel tellus semper finibus. In
  %       consequat, nibh sed mattis luctus, augue diam fermentum lectus.
  %     \item \textbf{In euismod erat metus} non ex. Vestibulum luctus augue in
  %       mi condimentum, at sollicitudin lorem viverra.
  %     \item \textbf{Suspendisse vulputate} mauris vel placerat consectetur.
  %       Mauris semper, purus ac hendrerit molestie, elit mi dignissim odio, in
  %       suscipit felis sapien vel ex.
  %   \end{itemize}

  %   Aenean tincidunt risus eros, at gravida lorem sagittis vel. Vestibulum ante
  %   ipsum primis in faucibus orci luctus et ultrices posuere cubilia Curae.

  % \end{alertblock}

    \begin{block}{Training Loop}

    \includegraphics[width=1\linewidth]{poster/resources/flow.png}

  \end{block}

  \end{column}

\separatorcolumn

\begin{column}{\colwidth}

  \begin{block}{Guess the Loss!}

  \vspace{0.5em}

  \textbf{Sequence A:} \\
  $1, 1, 2, 5, 12, 27, 58, 121, \dots$

  \begin{enumerate}
    \item $2^n - n$
    \item $a_n = a_{n-1} + a_{n-2}$
    \item \texttt{the pattern is [1, 1, 2, 5, 12, 27, 58, 121]}
  \end{enumerate}

  \vspace{0.75em}

  \textbf{Sequence B:} \\
  $6, 7, 13, 20, 33, 53, 86, 139, \dots$

  \begin{enumerate}
    \item $a_n = a_{n-1} + a_{n-2}$
    \item $a_n = 2n + 5$
    \item \texttt{the pattern is [6, 7, 13, 20, 33, 53, 86, 139]}
  \end{enumerate}

  \vspace{0.75em}

  % \textbf{Sequence C:} \\
  % $5, 8, 11, 14, 17, 20, 23, 26, \dots$

  % \begin{enumerate}
  %   \item $a_n = 3n + 5$
  %   \item $a_n = 2^n$
  %   \item \texttt{the pattern is [5, 8, 11, 14, 17, 20, 23, 26]}
  % \end{enumerate}

  % \vspace{0.5em}

  % \textit{Hint: MDL balances simplicity and predictive accuracy.}

\end{block}


\begin{block}{Results}

\small

We compare explanation scoring for the sequence
$(1, 1, 2, 5, 12, 27, 58, 121, 248, 503)$.

\begin{table}
\centering
\scriptsize
\scalebox{1.4}{
\begin{tabular}{l r r r}
\toprule
\textbf{Explanation $z$} & \textbf{Cond.} & \textbf{MDL} & \textbf{Prior / Likelihood} \\
\midrule
Base model output & 34.56 & 47.87 & 23.62 / 40.79 \\
Fine-tuned (LoRA) & \textbf{0.23} & \textbf{29.38} & 38.91 / 17.71 \\
\midrule
$2^n - n$ (true rule) & 39.35 & 29.38 & 38.91 / 17.71 \\
$3n + 1$ (false) & 42.27 & 93.71 & 34.47 / 59.24 \\
Dummy (copy data) & 63.88 & 95.77 & 95.77 / 0.00 \\
\bottomrule
\end{tabular}
}
\caption{Lower is better.}
\end{table}

\smallskip

\textbf{Key observation:}
\begin{itemize}
  \item Conditional likelihood favors what the model would say.
  \item MDL favors concise, predictive explanations.
\end{itemize}

\textbf{Training:}
We fine-tune a LoRA policy using REINFORCE with an MDL reward.
To improve exploration, we inject the correct rule with small
probability during sampling.

\smallskip

After fine-tuning, the model converges to the correct rule
$2^n - n$, achieving both low MDL loss and high confidence.

\end{block}

  \begin{block}{Conclusion}

  Fine-tuning a language model with an MDL-style reward
  (via LoRA) shifts it from predicting likely text to
  generating better explanations.

  Instead of optimizing only $p(z \mid x)$, the model is
  encouraged to balance:
  \begin{itemize}
    \item plausibility: $p(z)$
    \item correctness: $p(x \mid z)$
  \end{itemize}

  This leads to more concise, predictive, and interpretable
  rules, rather than fluent but incorrect explanations.

  The approach is domain-agnostic and can be applied beyond
  number sequences to general text-based pattern discovery.

\end{block}

  \begin{block}{References}

    \nocite{*}
    \footnotesize{\bibliographystyle{plain}\bibliography{poster}}

  \end{block}

\end{column}

\separatorcolumn
\end{columns}
\end{frame}

\end{document}
